% GENERATED FILE. Edit canonical lessons, then run npm run generate:books.
% canonical-source-hash: fnv1a64:10db1cf841c6e263
% canonical-lessons: MW-W35-digits-zero-one-two, MW-C35-read-ten-ticket, MW-W35-digits-three-four-five, MW-C35-read-five-ticket, MW-W35-digits-six-seven-eight, MW-C35-read-eight-ticket, MW-W35-digit-nine, MW-C35-read-ticket, MW-C35-ticket-four, MW-R35-script-close

\chapter{The Ten Digit Signs}
\label{ch:mw-digits}
\hlchaptermodality{35}

\begin{chapteropening}
\textbf{By the end of this chapter:} \emph{I can read a price written in Devanagari figures and write one from dictation.}

Ten digit signs arrive three at a time, each set cashed against the words it writes, and the four that a Latin eye reads wrongly are drilled apart from the six that it guesses correctly.
\end{chapteropening}

\section[\texorpdfstring{\mw{०} \mw{१} \mw{२} (0 1 2)}{० १ २ (0 1 2)}]{\mw{०} \mw{१} \mw{२} — the first three digits}

\label{lesson:MW-W35-digits-zero-one-two}



\begin{quote}
Recall the numbered-price payoff, then write ten, twenty and a hundred in words.
\end{quote}

\subsection*{Script — the first three figures}
Everything this chapter has written so far has been a number in words. A Rajasthani stall does not label anything in words. It writes figures, and Devanagari has its own ten.

\begin{quote}
\textbf{\mw{०}} — zero  \textbf{\mw{१}} — one  \textbf{\mw{२}} — two
\end{quote}

Look at where the ink sits. Each of these hangs from a stroke along the top, the same head-line that carries \textbf{\mw{क}} and \textbf{\mw{न}} and every other letter in this book. That is the quickest way to tell a Devanagari digit from a Latin one on a mixed sign.

\subsection*{Your turn}
Trace \textbf{\mw{०}}, \textbf{\mw{१}}, \textbf{\mw{२}} three times each. Say the Marwadi word for the second and third as you write them.

\begin{tcolorbox}[breakable,colback=teal!4,colframe=teal!35!black,title={Before you move on}]
Sources: \href{https://www.unicode.org/charts/PDF/U0900.pdf}{Unicode Devanagari chart}; \href{https://www.omniglot.com/language/numbers/rajasthani.htm}{Omniglot, Numbers in Rajasthani}.
\end{tcolorbox}
\section[Practice]{\mw{१०} — reading ten off a ticket}

\label{lesson:MW-C35-read-ten-ticket}



\begin{quote}
Say very. \emph{Write it:} the price question with its answer

\begin{itemize}
\raggedright
  \item \emph{Recall:} \emph{bīs}, and why twenty had to be taught before the numbers under it
\end{itemize}
\end{quote}

\subsection*{Your turn}
\begin{quote}
\textbf{\mw{१०}} — \emph{das} — ten
\end{quote}

The three signs from the last lesson already spell three amounts this book knows.

\begin{enumerate}
\raggedright
  \item \emph{Read it:} \textbf{\mw{१}}, \textbf{\mw{२}} and \textbf{\mw{१०}} aloud as \emph{ek}, \emph{do} and \emph{das}
  \item \emph{Write it:} each of the three in words, underneath its figure
  \item \emph{Cover:} the figures
  \item \emph{Write it:} the figures again, from the words
  \item Hear \emph{das}. \emph{Point to:} \textbf{\mw{१०}} rather than \textbf{\mw{दस}}
\end{enumerate}

Reading a figure and reading a word are two different skills, and a stall gives the learner the figure.

\begin{tcolorbox}[breakable,colback=teal!4,colframe=teal!35!black,title={Before you move on}]
Sources: \href{https://www.marwaripathshala.com/marwari-lesson-9-english}{Marwari Pathshala, Lesson 9}; \href{https://www.unicode.org/charts/PDF/U0900.pdf}{Unicode Devanagari chart}.
\end{tcolorbox}
\section[\texorpdfstring{\mw{३} \mw{४} \mw{५} (3 4 5)}{३ ४ ५ (3 4 5)}]{\mw{३} \mw{४} \mw{५} — three digits, two of them traps}

\label{lesson:MW-W35-digits-three-four-five}



\begin{quote}
Write very, write three in words, then write the digit for one.
\end{quote}

\subsection*{Script — three more figures}
\begin{quote}
\textbf{\mw{३}} — three  \textbf{\mw{४}} — four  \textbf{\mw{५}} — five
\end{quote}

Take them one at a time and set each beside the Latin figure a reader of this book already knows.

\textbf{\mw{३}} gives little trouble. \textbf{\mw{४}} and \textbf{\mw{५}} give a great deal: a reader who glances at \textbf{\mw{४}} sees an eight, and a reader who glances at \textbf{\mw{५}} sees a four. Neither guess is right, and both are wrong in the direction that matters at a counter.

Read them slowly for now. Speed on these two is worth nothing.

\subsection*{Your turn}
Trace \textbf{\mw{३}}, \textbf{\mw{४}}, \textbf{\mw{५}} three times each, saying \emph{tīn}, \emph{chār}, \emph{pāṅch} as you go. Then write \textbf{\mw{४}} and \textbf{\mw{५}} side by side and say aloud which is which without looking at this page.

\begin{tcolorbox}[breakable,colback=teal!4,colframe=teal!35!black,title={Before you move on}]
Sources: \href{https://www.unicode.org/charts/PDF/U0900.pdf}{Unicode Devanagari chart}; \href{https://www.omniglot.com/language/numbers/rajasthani.htm}{Omniglot, Numbers in Rajasthani}.
\end{tcolorbox}
\section[Practice]{\mw{३} \mw{४} \mw{५} against their words}

\label{lesson:MW-C35-read-five-ticket}



\begin{quote}
Say the too-expensive answer, and say an offer of five aloud. \emph{Write it:} an offer of ten

\begin{itemize}
\raggedright
  \item \emph{Recall:} \emph{so}, and the two other words that rhyme with it
\end{itemize}
\end{quote}

\subsection*{Your turn}
Run each digit in both directions. Reading only one way is how a lookalike survives practice.

\begin{enumerate}
\raggedright
  \item \emph{Read it:} \textbf{\mw{३}}, \textbf{\mw{५}}, \textbf{\mw{४}}, \textbf{\mw{१}}, \textbf{\mw{१०}} aloud, in that order
  \item \emph{Write it:} the word under each figure — \textbf{\mw{तीन}}, \textbf{\mw{पांच}}, \textbf{\mw{चार}}, \textbf{\mw{एक}}, \textbf{\mw{दस}}
  \item Now go the other way: hear \emph{chār} and \emph{pāṅch}. \emph{Write it:} the FIGURE, not the word
  \item \emph{Write it:} \textbf{\mw{४}} and \textbf{\mw{५}} from dictation, five times, given out of order
\end{enumerate}

\begin{tcolorbox}[breakable,colback=teal!4,colframe=teal!35!black,title={Before you move on}]
Sources: \href{https://www.marwaripathshala.com/marwari-lesson-9-english}{Marwari Pathshala, Lesson 9}; \href{https://www.unicode.org/charts/PDF/U0900.pdf}{Unicode Devanagari chart}.
\end{tcolorbox}
\section[\texorpdfstring{\mw{६} \mw{७} \mw{८} (6 7 8)}{६ ७ ८ (6 7 8)}]{\mw{६} \mw{७} \mw{८} — the last two traps}

\label{lesson:MW-W35-digits-six-seven-eight}



\begin{quote}
Write the too-expensive answer, then write the digits for four and five.
\end{quote}

\subsection*{Script — the last three of the traps}
\begin{quote}
\textbf{\mw{६}} — six  \textbf{\mw{७}} — seven  \textbf{\mw{८}} — eight
\end{quote}

\textbf{\mw{६}} is close enough to its Latin cousin that a reader guesses it correctly. The other two are not.

\textbf{\mw{७}} carries a hook that pulls the eye towards a nine. \textbf{\mw{८}} is closed and round and looks more like a letter than a figure. Both stand next to \textbf{\mw{६}} here on purpose: a trap is easiest to learn beside something safe.

\subsection*{Your turn}
Trace \textbf{\mw{६}}, \textbf{\mw{७}}, \textbf{\mw{८}} three times each, saying \emph{chha}, \emph{sāt}, \emph{āṭh} as you go. Then cover the models and write \textbf{\mw{७}} and \textbf{\mw{८}} from the words alone.

\begin{tcolorbox}[breakable,colback=teal!4,colframe=teal!35!black,title={Before you move on}]
Sources: \href{https://www.unicode.org/charts/PDF/U0900.pdf}{Unicode Devanagari chart}; \href{https://www.omniglot.com/language/numbers/rajasthani.htm}{Omniglot, Numbers in Rajasthani}.
\end{tcolorbox}
\section[Practice]{\mw{८} and \mw{२०} — one figure, then two}

\label{lesson:MW-C35-read-eight-ticket}



\begin{quote}
\emph{Write it:} the price question with its answer, then the digits for six and seven

\begin{itemize}
\raggedright
  \item \emph{Recall:} the danda that closes a one-word price, two chapters back
\end{itemize}
\end{quote}

\subsection*{Your turn}
\begin{quote}
\textbf{\mw{८}} — \emph{āṭh} — eight
\end{quote}

\begin{quote}
\textbf{\mw{२०}} — \emph{bīs} — twenty
\end{quote}

\textbf{\mw{२०}} is not a two and a zero read one after the other. It is one amount, and its name is a word with no two and no ten inside it.

\begin{enumerate}
\raggedright
  \item \emph{Read it:} \textbf{\mw{६}}, \textbf{\mw{७}}, \textbf{\mw{८}} aloud, then \textbf{\mw{१०}} and \textbf{\mw{२०}}
  \item \emph{Write it:} the word under each of the five
  \item Hear \emph{sāt}, \emph{chha}, \emph{bīs}. \emph{Write it:} the FIGURE each time
  \item \emph{Write it:} \textbf{\mw{२०}} and \textbf{\mw{८}} from dictation, given out of order
\end{enumerate}

\begin{tcolorbox}[breakable,colback=teal!4,colframe=teal!35!black,title={Before you move on}]
Sources: \href{https://www.marwaripathshala.com/marwari-lesson-9-english}{Marwari Pathshala, Lesson 9}; \href{https://www.unicode.org/charts/PDF/U0900.pdf}{Unicode Devanagari chart}.
\end{tcolorbox}
\section[\texorpdfstring{\mw{९} (9)}{९ (9)}]{\mw{९} — the tenth digit}

\label{lesson:MW-W35-digit-nine}



\begin{quote}
Say a little, then write the digits for seven and eight.
\end{quote}

\subsection*{Script — the tenth figure}
\begin{quote}
\textbf{\mw{९}} — nine — \emph{no}
\end{quote}

Like \textbf{\mw{३}} and \textbf{\mw{६}}, this one a Latin eye can very nearly guess, and unlike \textbf{\mw{७}} the guess is right.

With \textbf{\mw{९}} the set is closed. Ten digit signs, taught across four lessons, and the eye can now read any figure a stall prints as far as the words go.

Write \textbf{\mw{९}} and \textbf{\mw{नो}} side by side. The figure and the word for nine share nothing at all — which is the whole reason the figures had to be taught separately.

\subsection*{Your turn}
Trace \textbf{\mw{९}} three times. Then write the whole run \textbf{\mw{०} \mw{१} \mw{२} \mw{३} \mw{४} \mw{५} \mw{६} \mw{७} \mw{८} \mw{९}} once, slowly, saying the Marwadi word for each from one onward.

\begin{tcolorbox}[breakable,colback=teal!4,colframe=teal!35!black,title={Before you move on}]
Sources: \href{https://www.unicode.org/charts/PDF/U0900.pdf}{Unicode Devanagari chart}; \href{https://www.omniglot.com/language/numbers/rajasthani.htm}{Omniglot, Numbers in Rajasthani}.
\end{tcolorbox}
\section[Practice]{\mw{९}, \mw{२०}, \mw{१००} — one figure, two, three}

\label{lesson:MW-C35-read-ticket}



\begin{quote}
\emph{Write it:} the word for a little, then an offer of five

\begin{itemize}
\raggedright
  \item \emph{Recall:} \textbf{\mw{४}} and \textbf{\mw{५}}, and which Latin figures each is mistaken for
\end{itemize}
\end{quote}

\subsection*{Your turn}
\begin{quote}
\textbf{\mw{९}} — \emph{no}  \textbf{\mw{२०}} — \emph{bīs}  \textbf{\mw{१००}} — \emph{so}
\end{quote}

Three tickets, one to three figures long, and the words behind them run the other way: the longest ticket carries the shortest word.

\begin{enumerate}
\raggedright
  \item \emph{Read it:} the three aloud, then read them again shuffled
  \item \emph{Write it:} each amount in words, beneath its figures
  \item Hear \emph{so}, \emph{no}, \emph{bīs}. \emph{Write it:} the FIGURES each time
  \item \emph{Read it:} six mixed tickets — \textbf{\mw{५}}, \textbf{\mw{१००}}, \textbf{\mw{७}}, \textbf{\mw{२०}}, \textbf{\mw{३}}, \textbf{\mw{९}} — then say each amount without pausing on the trap digits
\end{enumerate}

\begin{tcolorbox}[breakable,colback=teal!4,colframe=teal!35!black,title={Before you move on}]
Sources: \href{https://www.marwaripathshala.com/marwari-lesson-9-english}{Marwari Pathshala, Lesson 9}; \href{https://www.unicode.org/charts/PDF/U0900.pdf}{Unicode Devanagari chart}.
\end{tcolorbox}
\section[Practice]{Payoff — the printed price}

\label{lesson:MW-C35-ticket-four}



\begin{quote}
Recall the numbered-price payoff, then say the make-it word.
\end{quote}

\subsection*{Your turn}
Four skills, scored separately.

\begin{enumerate}
\raggedright
  \item \textbf{Listening.} Hear twelve amounts named. \emph{Write it:} each as a Devanagari figure
  \item \textbf{Speaking.} \emph{Read it:} twelve printed tickets aloud as Marwadi words, including every one of \textbf{\mw{४}}, \textbf{\mw{५}}, \textbf{\mw{७}} and \textbf{\mw{८}}
  \item \emph{Read it:} the labels on a stall of six goods, then say which is dearest and which is cheapest
  \item \textbf{Writing.} \emph{Write it:} all ten digits from dictation, out of order, then \textbf{\mw{१०}}, \textbf{\mw{२०}} and \textbf{\mw{१००}}
\end{enumerate}

Pass each separately. A learner who reads the safe digits fluently and stalls on the four traps has not passed; those four are the ones a ticket will punish.

\begin{tcolorbox}[breakable,colback=teal!4,colframe=teal!35!black,title={Before you move on}]
Sources: \href{https://www.marwaripathshala.com/marwari-lesson-9-english}{Marwari Pathshala, Lesson 9}; \href{https://www.unicode.org/charts/PDF/U0900.pdf}{Unicode Devanagari chart}.
\end{tcolorbox}
\section[Practice]{Close the writing loop for \mw{४}, \mw{५}, \mw{७} and \mw{८}}

\label{lesson:MW-R35-script-close}



\begin{quote}
Recall the ticket payoff. \emph{Write it:} the make-it word and an offer of twenty
\end{quote}

\subsection*{Your turn}
\emph{Write it:} four single figures from spoken cues, given out of order, twice through — \textbf{\mw{४}}, \textbf{\mw{५}}, \textbf{\mw{७}}, \textbf{\mw{८}}

\emph{Write it:} the word beside each — \textbf{\mw{चार}}, \textbf{\mw{पांच}}, \textbf{\mw{सात}}, \textbf{\mw{आठ}}; repair only the missed figure and write its pair once more

\begin{tcolorbox}[breakable,colback=teal!4,colframe=teal!35!black,title={Before you move on}]
Sources: \href{https://www.unicode.org/charts/PDF/U0900.pdf}{Unicode Devanagari chart}; \href{https://www.omniglot.com/language/numbers/rajasthani.htm}{Omniglot, Numbers in Rajasthani}.
\end{tcolorbox}
